{\color{red}
\section{Sugestão de uso em sala}\label{sec:roteiro}

O material foi pensado para disciplinas de laboratório, de tratamento de dados ou de métodos computacionais dos anos finais da graduação em física e em engenharia, e pode ser usado também no início da pós-graduação. Supõe-se que o estudante conheça probabilidade básica (função distribuição acumulada, densidade, média e variância), cálculo de uma variável (derivada, integral e regra da cadeia) e o suficiente de Python para executar e modificar as células de um notebook. Uma sequência possível ocupa duas aulas de cerca de 100 minutos. Na primeira, trabalham-se as seções~\ref{sec:gld} e \ref{sec:dados} com a figura~\ref{fig:parametros} e o Código~\ref{lst:quantil}. Na segunda, o estudo do vento ou o da temperatura é refeito no notebook, e o outro fica como tarefa extraclasse. A duração é uma estimativa dos autores, ainda não verificada em turma.

A tabela~\ref{tab:roteiro} liga cada objetivo da seção~\ref{sec:introducao} a uma atividade e ao que se espera que o estudante produza. Em todas as tarefas, recomenda-se pedir primeiro uma previsão por escrito e só depois a execução do código, para que a discussão parta da diferença entre o previsto e o obtido.

\begin{table}[htbp]
\centering
\caption{Objetivos, atividades e evidências esperadas na sequência proposta.}
\label{tab:roteiro}
\small
\begin{tabular}{p{3.6cm} p{5.4cm} p{5.0cm}}
\toprule
\textbf{Objetivo} & \textbf{Atividade} & \textbf{Evidência esperada} \\
\midrule
Interpretar a função quantil & Prever o efeito de dobrar $\lambda_2$ e conferir no notebook & Explicação de que a dispersão cai à metade e a forma não muda \\
Compreender o ajuste por momentos & Comparar os momentos e o suporte da GLD ajustada ao vento & Constatação de que momentos iguais não garantem suporte adequado \\
Avaliar criticamente um ajuste & Comparar $\Delta_p$ da normal com a variabilidade dos quantis de julho & Justificativa escrita para manter, ou não, a normal \\
Reconhecer limites da interpretação física & Discutir a potência eólica a partir de médias diárias & Argumento de que médias diárias não determinam $\langle v^3\rangle$ \\
\bottomrule
\end{tabular}
\end{table}

As quatro tarefas abaixo correspondem às linhas da tabela, com as respostas esperadas.

\begin{enumerate}
    \item \textbf{Dobrar $\lambda_2$.} Pela equação~\eqref{eq:fkml}, $Q(u) - \lambda_1$ é proporcional a $1/\lambda_2$. Dobrar $\lambda_2$ divide por dois a distância de cada quantil a $\lambda_1$, e portanto o desvio padrão e a distância interquartil, sem alterar a assimetria nem a curtose. O painel (b) da figura~\ref{fig:parametros} permite conferir.
    \item \textbf{Momentos e suporte.} Os momentos são integrais sobre toda a distribuição e podem coincidir com os da amostra mesmo que o limite inferior $a = \lambda_1 - 1/(\lambda_2\lambda_3)$ fique acima do menor valor observado. Espera-se que o estudante calcule $a = 1,148\,\text{m/s}$, conte os 84 dias abaixo dele e explique por que a verossimilhança se anula.
    \item \textbf{Manter a normal em julho?} Os desvios da normal ($+0,16$, $-0,09$ e $+0,12\,^\circ$C) são menores que os desvios padrão por reamostragem de anos ($0,38$, $0,10$ e $0,21\,^\circ$C). Uma resposta adequada mantém a normal para descrever as temperaturas típicas e reconhece que, se o interesse for a cauda fria, a assimetria capturada pela GLD pode importar.
    \item \textbf{Médias diárias e potência eólica.} A densidade de potência depende de $\langle v^3\rangle$ calculado com velocidades em intervalos curtos. Como $v^3$ é convexa, a média de $v^3$ ao longo de um dia é maior ou igual ao cubo da média diária, de modo que usar médias diárias subestima a potência. A resposta esperada é que seriam necessários dados horários ou de dez minutos.
\end{enumerate}
\par}
